\title{Group Sequence Policy Optimization}

\begin{abstract}
We propose Group Sequence Policy Optimization (GSPO), a robust and highly effective reinforcement learning method tailored for the training of large language models. In contrast to earlier methods that focus on token-wise importance weights, GSPO computes importance ratios over entire sequences, enabling sequence-based clipping, rewarding, and policy updates. Our experiments show that GSPO offers notable gains in both training speed and final performance over the GRPO method, brings increased stability to Mixture-of-Experts (MoE) RL training, and offers promise in streamlining RL system design. These advantages have played a critical role in elevating the performance of recent Qwen3 models.
\end{abstract}

\section{Introduction}
\label{sec:intro}

Reinforcement learning (RL) has become a central approach for advancing the capabilities of language models \citep{o1,dpsk-r1,qwq32b,qwen3}. When scaled to large training regimes, RL enables language models to address complex challenges, including those found in competitive mathematics and programming, by supporting more advanced and extended reasoning abilities.

For RL to effectively harness increased computational resources, it is essential to ensure that the training process remains stable and resilient. Nonetheless, leading RL methods such as GRPO \citep{grpo} are plagued by significant instability during the training of large-scale language models, frequently leading to disastrous and sometimes permanent model failures \citep{qwen3, minimax-m1}. Such instability poses a substantial barrier to extending the performance and reasoning depth of language models using RL techniques.

This work demonstrates that the core instability observed in GRPO originates from a fundamental misuse and breakdown of importance sampling weights within the framework of its algorithm. Such misapplication leads to pronounced variance in training signals, with the resultant noise compounding as response sequences become longer. Furthermore, this effect is exacerbated by the application of the clipping procedure, which collectively drives the model toward collapse.

In response to these inherent shortcomings, we introduce \textbf{Group Sequence Policy Optimization (GSPO)}, a novel reinforcement learning technique tailored for the training of large language models. GSPO’s central contribution is its rigorous and principled formulation of the importance ratio, grounded in sequence likelihood as detailed by \citet{click}, adhering closely to the foundational tenets of importance sampling. Moreover, GSPO evaluates normalized rewards by considering the advantages accrued from multiple responses per query, thus preserving the coherence between the process of rewarding at the sequence level and subsequent optimization steps.

Our experimental results reveal that GSPO markedly outperforms GRPO in terms of training stability, efficiency, and overall effectiveness. Notably, GSPO naturally addresses stability issues encountered in RL training for large Mixture-of-Experts (MoE) architectures, thus removing the reliance on intricate stabilization techniques and paving the way for a more streamlined RL infrastructure. These strengths of GSPO have played a pivotal role in the remarkable performance gains observed in the most recent Qwen3 models. We foresee GSPO serving as a resilient and scalable algorithmic base, facilitating ongoing progress in the large-scale RL training of language models.

\section{Preliminaries}

\paragraph{Notation}
Throughout this work, we represent an autoregressive language model with parameters $\theta$ as a policy $\pi_\theta$. A query is denoted by $x$, while the collection of all queries is represented by $\mathcal{D}$. For any response $y$ generated in reply to $x$, the probability of $y$ under the policy $\pi_\theta$ is written as $\pi_\theta (y | x)=\prod_{t=1}^{|y|} \pi_\theta (y_t | x, y_{<t} )$, where $|y|$ indicates the total token count in $y$. Each query-response tuple $(x, y)$ is assigned a reward $r(x, y) \in [ 0, 1 ]$ by an external verifier $r$.

\paragraph{Proximal Policy Optimization (PPO)} Using samples generated from the old policy $\pi_{\theta_\text{old}}$, PPO \citep{ppo} constrains the policy update within a proximal region of the old policy through the clipping mechanism. Specifically, PPO employs the following objective for policy optimization (we omit the KL regularization term hereinafter for brevity, as it is not the focus of this paper):

\begin{align}
\mathcal{J}_\text{PPO}(\theta) = \mathbb{E}_{ x \sim \mathcal{D},\, y \sim \pi_{\theta_\text{old}}( \cdot | x) }
\left[ \frac{1}{|y|} \sum_{t=1}^{|y|} 
\min \left( w_{t}(\theta) \widehat{A}_{t},  \, \mathrm{clip} \left( w_{t}(\theta), 1 - {\varepsilon}, 1 + {\varepsilon}\right) \widehat{A}_{t} \right)
\right],
\end{align}

In this context, the importance ratio for token $y_t$ is given by $
w_{t}(\theta) = \frac{ \pi_{\theta} (y_{t} | x, y_{<t}) }{ \pi_{\theta_\text{old}} (y_{t} | x,y_{<t})} $, where the advantage $\widehat{A}_{t}$ associated with $y_t$ is approximated using a separate value model, and $\varepsilon$ denotes the clipping threshold applied to the importance ratios.

One of the main difficulties associated with PPO is its significant dependence on the value model. Typically, this value model is comparable in complexity and size to the policy model itself, resulting in substantial overhead in terms of both memory and computation. The performance of the algorithm is further contingent upon the accuracy of the value model’s estimations. Building an effective value model is already non-trivial, and scaling such a model to accommodate longer output sequences or more intricate problems becomes increasingly problematic.

\paragraph{Group Relative Policy Optimization (GRPO)} GRPO \citep{grpo} bypasses the need for the value model by computing the relative advantage of each response within a group of responses to the same query. Specifically, GRPO optimizes the following objective:

\begin{align}
\label{equ:grpo}
\mathcal{J}_\text{GRPO}(\theta) = \mathbb{E}_{ x \sim \mathcal{D},\, \{y_i\}_{i=1}^G \sim \pi_{\theta_\text{old}}( \cdot | x) }
\left[ \frac{1}{G} \sum_{i=1}^{G} \frac{1}{|y_i|} \sum_{t=1}^{|y_i|} 
\min \left( w_{i,t}(\theta) \widehat{A}_{i,t},  \, \mathrm{clip} \left( w_{i,t}(\theta), 1 - {\varepsilon}, 1 + {\varepsilon}\right) \widehat{A}_{i,t} \right)
\right],
\end{align}

where $G$ is the number of generated responses to each query $x$ (i.e., the group size), and the importance ratio $w_{i,t}(\theta)$ and advantage $\widehat{A}_{i,t}$ of token $y_{i,t}$ are:

\begin{align}
    w_{i,t}(\theta)=\frac{ \pi_{\theta} (y_{i,t} | x, y_{i,<t}) }{ \pi_{\theta_\text{old}} (y_{i,t} | x,y_{i,<t})},\quad\ 
    \widehat{A}_{i,t} = \widehat{A}_{i} = \frac{r(x, y_i) - \mathrm{mean} \left( \{ r(x, y_i) \}_{i=1}^G \right) }{ \mathrm{std} \left( \{ r(x, y_i) \}_{i=1}^G \right) },
\end{align}

respectively, where all the tokens in $y_i$ share the same advantage as $\widehat{A}_{i}$.

\section{Motivation}
\label{sec:motivation}

As model sizes increase, alongside factors such as sparsity—common in Mixture-of-Experts architectures—and longer response generations, there arises a requirement for utilizing large rollout batch sizes to ensure effective hardware throughput during reinforcement learning. To enhance how efficiently data is used, it is conventional to divide these extensive rollout batches into smaller mini-batches, each contributing to separate gradient update steps. However, this batching strategy creates an inherent off-policy scenario: the outputs $y$ are typically produced under a previous policy $\pi_{\theta_\text{old}}$, which may differ from the currently optimized policy $\pi_\theta$. This situation underlies the rationale for incorporating clipping techniques in algorithms like PPO and GRPO, as these mechanisms restrict the influence of data points that are too "off-policy" and thus maintain the reliability of the gradient computation.

While mechanisms like clipping aim to manage this off-policy discrepancy, we identify a more fundamental issue in GRPO: \textit{its objective is ill-posed}. This problem becomes particularly acute when training large models on long-response tasks, leading to catastrophic model collapse. The ill-posed nature of the GRPO objective stems from a misapplication of importance sampling weights. The principle of importance sampling is to estimate the expectation of a function $f$ under a target distribution $\pi_\text{tar}$ by re-weighting samples drawn from a behavior distribution $\pi_\text{beh}$:

\begin{align}
\mathbb{E}_{z \sim \pi_\text{tar}} \left[ f(z) \right] = \mathbb{E}_{z \sim \pi_\text{beh}} \left[ \frac{ \pi_\text{tar} (z) }{ \pi_\text{beh} (z)} \, f(z) \right].
\end{align}

A key aspect of this method is the requirement to average over a large number of samples ($N \gg 1$) drawn from the behavior policy $\pi_\text{beh}$; only in this regime does the importance weight $\frac{ \pi_\text{tar} (z) }{ \pi_\text{beh} (z)}$ serve its intended function of ameliorating the mismatch between distributions.

By contrast, GRPO implements the importance weighting term $\frac{ \pi_{\theta} (y_{i,t} | x, y_{i,<t}) }{ \pi_{\theta_\text{old}} (y_{i,t} | x,y_{i,<t})}$ at each token timestep $t$. However, this token-wise weight is calculated on the basis of a solitary realization $y_{i,t}$ sampled from the next-token distribution $\pi_{\theta_\text{old}}( \cdot | x, y_{i, <t})$, thereby failing to adequately perform the desired correction for distributional shift. The result is the injection of substantial variance into the training gradients, which accumulates over extended output sequences and is further aggravated by the application of the clipping procedure. Our empirical findings show that such a combination can precipitate a catastrophic model collapse that, once set in, tends to be irreversible. After collapse, resuming successful training proves elusive—even reverting to earlier checkpoints, meticulously adjusting hyperparameter settings (such as the clipping bounds), lengthening generations, or substituting RL queries offer little remediation.

These empirical results reveal a deeper flaw in the architecture of GRPO. The ineffectiveness of token-level importance weights highlights a foundational insight: \textbf{the unit of optimization should correspond to the unit of reward assignment}. Since the reward signal pertains to the full output sequence, implementing off-policy correction at the granularity of individual tokens appears fundamentally flawed. This observation motivates us to abandon token-level objectives and, instead, directly leverage importance weights and perform optimization at the \textit{sequence level}.

\section{Algorithm}

\subsection{GSPO: Group Sequence Policy Optimization}

Although the token-level importance weight $\frac{ \pi_{\theta} (y_{i,t} | x, y_{i,<t}) }{ \pi_{\theta_\text{old}} (y_{i,t} | x,y_{i,<t})}$ presents challenges in GRPO, we find that, for language generation, the \textit{sequence-level} importance weight $\frac{ \pi_{\theta} (y | x) }{ \pi_{\theta_\text{old}} (y | x)}$ offers a well-defined theoretical interpretation. Specifically, it quantifies the divergence between a response $y$ generated by $\pi_{\theta_\text{old}} (\cdot | x)$ and the distribution $\pi_{\theta} (\cdot | x)$. This measurement intuitively matches sequence-level reward computation and effectively guides the use of a clipping strategy.

Based on this straightforward observation, we propose the \textbf{Group Sequence Policy Optimization (GSPO)} algorithm. GSPO employs the following sequence-level optimization objective:

\begin{align}
\mathcal{J}_\text{GSPO} (\theta) =
\mathbb{E}_{ x \sim \mathcal{D},\, \{y_i\}_{i=1}^G \sim \pi_{\theta_\text{old}}( \cdot | x) }
\left[ 
\frac{1}{G} \sum_{i=1}^{G}
\min \left( s_{i}(\theta)  \widehat{A}_{i},  \, \mathrm{clip} \left( s_{i}(\theta), 1 - {\varepsilon}, 1 + {\varepsilon} \right) \widehat{A}_{i} \right) 
\right],
\label{equ:gspo}
\end{align}

where we adopt the group-based advantage estimation:

\begin{align}
\widehat{A}_{i} = \frac{r(x, y_i) - \mathrm{mean} \left( \{ r(x, y_i) \}_{i=1}^G \right) }{ \mathrm{std} \left( \{ r(x, y_i) \}_{i=1}^G \right) },
\end{align}

and define the importance ratio $s_{i}(\theta)$ based on sequence likelihood \citep{click}:

\begin{align}
s_{i}(\theta) = \left( \frac{ \pi_{\theta} (y_i | x) }{ \pi_{\theta_\text{old}} (y_i | x)} \right)^{\frac{1}{|y_i|}}
=
\exp \left( \frac{1}{|y_i|} \sum_{t=1}^{|y_i|} \log \frac{ \pi_{\theta} (y_{i,t} | x, y_{i,<t}) }{ \pi_{\theta_\text{old}} (y_{i,t} | x,y_{i,<t})} \right).
\end{align}

Consequently, GSPO implements response-level clipping rather than token-level clipping, effectively removing highly ``off-policy'' examples from the gradient computation. This approach is consistent with both sequence-level reward assignment and the optimization process. Additionally, we apply length normalization to $s_{i}(\theta)$ in order to limit its variability and ensure that $s_{i}(\theta)$ stays within a consistent numeric interval. Without normalization, changes in the probabilities of a handful of tokens could cause excessive volatility in the overall importance ratio for the sequence, and responses of differing lengths would necessitate different clipping thresholds. It is worth mentioning that the scale of clipping in GSPO generally diverges—by at least an order of magnitude—from that used in earlier approaches such as GRPO, owing to the unique ways in which importance ratios are computed.

\subsection{Gradient Analysis}
\label{subsec:gradient}

We can derive the gradient of the GSPO objective as follows (clipping is omitted for brevity):

\begin{align}
\nabla_{\theta} \mathcal{J}_\text{GSPO} (\theta)
=&\ 
\nabla_{\theta} \mathbb{E}_{ x \sim \mathcal{D},\, \{y_i\}_{i=1}^G \sim \pi_{\theta_\text{old}}( \cdot | x) }
\left[ 
\frac{1}{G} \sum_{i=1}^{G} s_{i}(\theta) \widehat{A}_{i}
\right] \\
= &\ 
\mathbb{E}_{ x \sim \mathcal{D},\, \{y_i\}_{i=1}^G \sim \pi_{\theta_\text{old}}( \cdot | x) }
\left[ 
\frac{1}{G} \sum_{i=1}^{G}
s_{i}(\theta) \widehat{A}_{i}
\cdot \nabla_{\theta} \log s_{i}(\theta)
\right] \\
=&\ 
\mathbb{E}_{ x \sim \mathcal{D},\, \{y_i\}_{i=1}^G \sim \pi_{\theta_\text{old}}( \cdot | x) }
\left[ 
\frac{1}{G} \sum_{i=1}^{G} 
\left( \frac{ \pi_{\theta} (y_i | x) }{ \pi_{\theta_\text{old}} (y_i | x)} \right)^{\frac{1}{|y_i|}} \widehat{A}_{i} 
\cdot \frac{1}{|y_i|} \sum_{t=1}^{|y_i|} \nabla_{\theta} \log \pi_{\theta} (y_{i,t} | x, y_{i,<t}) 
\right].
\label{equ:gspo_grad}
\end{align}

For comparison, the gradient of the GRPO objective is as follows (note that $\widehat{A}_{i,t} = \widehat{A}_{i}$):

\begin{align}
\nabla_{\theta} \mathcal{J}_\text{GRPO} (\theta) 
=&\ 
\nabla_{\theta} \mathbb{E}_{ x \sim \mathcal{D},\, \{y_i\}_{i=1}^G \sim \pi_{\theta_\text{old}}( \cdot | x) }
\left[ \frac{1}{G} \sum_{i=1}^{G} \frac{1}{|y_i|} \sum_{t=1}^{|y_i|} 
w_{i,t}(\theta) \widehat{A}_{i,t}
\right] \\
=&\
\mathbb{E}_{ x \sim \mathcal{D},\, \{y_i\}_{i=1}^G \sim \pi_{\theta_\text{old}}( \cdot | x) }
\left[ \frac{1}{G} \sum_{i=1}^{G} \widehat{A}_{i} 
\cdot \frac{1}{|y_i|} \sum_{t=1}^{|y_i|} 
\frac{ \pi_{\theta} (y_{i,t} | x, y_{i,<t}) }{ \pi_{\theta_\text{old}} (y_{i,t} | x,y_{i,<t})} 
\nabla_{\theta} \log \pi_{\theta} (y_{i,t} | x, y_{i,<t})  
\right].
\end{align}

The key difference between GSPO and GRPO centers on \textit{the manner in which gradients of the token log likelihoods are weighted}. In the case of GRPO, each token receives a weight based on its individual ``importance weight'' $\frac{ \pi_{\theta} (y_{i,t} | x, y_{i,<t}) }{ \pi_{\theta_\text{old}} (y_{i,t} | x,y_{i,<t})}$, resulting in varying token contributions. These weights may range over $(0, 1+\varepsilon]$ when $\widehat{A}_{i} >0$ or $[1-\varepsilon, +\infty)$ for $\widehat{A}_{i} < 0$, and their variability is significant; as training continues, their cumulative effects can introduce instability. On the other hand, GSPO applies uniform weighting across all tokens within a response, thereby mitigating the instability present in GRPO.

\subsection{GSPO-token: A Token-level Objective Variant}

In scenarios like multi-turn RL, we may desire a finer-grained advantage adjustment than the sequence level. To this end, we introduce a token-level objective variant of GSPO, namely \textbf{GSPO-token}, to allow token-wise advantage customization:

\begin{align}
\mathcal{J}_\text{GSPO-token}(\theta)
=
\mathbb{E}_{ x \sim \mathcal{D},\, \{y_i\}_{i=1}^G \sim \pi_{\theta_\text{old}}( \cdot | x) }
\left[ 
\frac{1}{G} \sum_{i=1}^{G} \frac{1}{|y_i|} \sum_{t=1}^{|y_i|} 
\min \left( s_{i,t}(\theta) \widehat{A}_{i,t},  \, \mathrm{clip} \left( s_{i,t}(\theta), 1 - {\varepsilon}, 1 + {\varepsilon} \right) \widehat{A}_{i,t} \right) 
\right],
\label{equ:gspo-token}
\end{align}

where

\begin{align}
s_{i,t}(\theta) 
= \mathrm{sg} \left[ s_{i}(\theta) \right]  \cdot \frac{ \pi_{\theta} (y_{i,t} | x, y_{i,<t}) }{ \mathrm{sg} \left[ \pi_{\theta} (y_{i,t} | x, y_{i,<t}) \right] },
\end{align}

and $\mathrm{sg}[\cdot]$ denotes only taking the numerical value but stopping the gradient, corresponding to the \texttt{detach} operation in PyTorch. The gradient of GSPO-token can be derived as:

\begin{align}
\nabla_{\theta} \mathcal{J}_\text{GSPO-token}(\theta)
=&\ 
\nabla_{\theta} \mathbb{E}_{ x \sim \mathcal{D},\, \{y_i\}_{i=1}^G \sim \pi_{\theta_\text{old}}( \cdot | x) }
\left[ 
\frac{1}{G} \sum_{i=1}^{G}
\frac{1}{|y_i|} \sum_{t=1}^{|y_i|} 
s_{i,t}(\theta) \widehat{A}_{i,t}
\right] \\
=&\ 
\mathbb{E}_{ x \sim \mathcal{D},\, \{y_i\}_{i=1}^G \sim \pi_{\theta_\text{old}}( \cdot | x) }
\left[ 
\frac{1}{G} \sum_{i=1}^{G} s_i (\theta) \cdot \frac{1}{|y_i|} \sum_{t=1}^{|y_i|} 
\widehat{A}_{i,t} \frac{ \nabla_{\theta} \pi_{\theta} (y_{i,t} | x, y_{i,<t}) }{ \pi_{\theta} (y_{i,t} | x, y_{i,<t}) }
\right] \\
=&\ 
\mathbb{E}_{ x \sim \mathcal{D},\, \{y_i\}_{i=1}^G \sim \pi_{\theta_\text{old}}( \cdot | x) }
\left[ 
\frac{1}{G} \sum_{i=1}^{G}  \left( \frac{ \pi_{\theta} (y_i | x) }{ \pi_{\theta_\text{old}} (y_i | x)} \right)^{\frac{1}{|y_i|}}
\cdot \frac{1}{|y_i|} \sum_{t=1}^{|y_i|}
\widehat{A}_{i,t} \nabla_{\theta} \log \pi_{\theta} (y_{i,t} | x, y_{i,<t})  
\right].
\label{equ:gspo-token_grad}
\end{align}

Observe that the expression $\frac{ \pi_{\theta} (y_{i,t} | x, y_{i,<t}) }{ \mathrm{sg} \left[ \pi_{\theta} (y_{i,t} | x, y_{i,<t}) \right] }$ always evaluates to 1, implying that $s_{i,t}(\theta)$ has the same numerical value as $s_{i}(\theta)$. A comparison between Equation~\eqref{equ:gspo} and~\eqref{equ:gspo-token}, as well as between Equation~\eqref{equ:gspo_grad} and~\eqref{equ:gspo-token_grad}, reveals that GSPO-token and GSPO yield numerically equivalent results regarding their objective functions, clipping conditions, and theoretical gradients, provided all token-wise advantages in response $y_i$ are set uniformly, namely $\widehat{A}_{i,t} = \widehat{A}_{i}$. Nevertheless, GSPO-token provides the added advantage of permitting token-level customization of advantage values.

\section{Experiments and Discussion}

\subsection{Empirical Results}

We conduct experiments utilizing a cold-start model that has been fine-tuned from Qwen3-30B-A3B-Base, presenting both the progression of training rewards and the resulting performance on several benchmarks, including AIME'24 (using the mean Pass@1 across 32 samples), LiveCodeBench (from 202410-202502, with the mean Pass@1 based on 8 samples), and CodeForces (measured by Elo Rating). During reinforcement learning training, each rollout batch is subdivided into four mini-batches for the purpose of gradient updating. For GSPO, the left and right clipping boundaries in Equation~\eqref{equ:gspo} are configured at 3e-4 and 4e-4, correspondingly. To provide a point of comparison, we employ GRPO as the baseline, setting the clipping intervals in Equation~\eqref{equ:grpo} to 0.2 and 0.27; these values have been systematically tuned to ensure equitable evaluation. It is important to highlight that GRPO requires the Routing Replay approach to facilitate the stable convergence of MoE RL—this aspect will be addressed in \S~\ref{subsec:routing_replay}—while \textit{GSPO eliminates the requirement for such a strategy}.

As illustrated in Figure~\ref{fig:results}, training using GSPO maintains stability over the entire process. Notably, \textbf{GSPO enables ongoing gains in performance by scaling up training compute, frequently refreshing the query set, and increasing the generation length}. In addition, GSPO exhibits enhanced training efficiency compared to GRPO, as evidenced by superior training accuracy and benchmark results when constrained to identical levels of training compute and query consumption. Furthermore, we have deployed GSPO in the RL training of the most recent Qwen3 models, providing strong empirical support for GSPO's capacity to harness RL scaling to improve large language model performance.

\subsection{Curious Observation on Clipping Fractions}

An important difference between GSPO and GRPO lies in GSPO's method of clipping entire sequences as opposed to clipping individual tokens. As illustrated in Figure~\ref{fig:clipping}, there is an observed gap of nearly two orders of magnitude between GSPO and GRPO regarding the proportion of tokens being clipped, and this remains consistent irrespective of how the clipping thresholds are set. Interestingly, GSPO manages to maintain greater training efficiency than GRPO, even though it excludes a substantially larger number of tokens (thus reducing the number used for training or estimating gradients). This somewhat surprising result—that discarding a much higher share of tokens can still yield superior efficiency—suggests that the gradient estimates obtained at the token level in GRPO are notably noisy and suboptimal for effective use of samples. Conversely, by utilizing gradients at the sequence level, GSPO is able to deliver a stronger and more useful training signal.

\subsection{Benefit of GSPO for MoE Training}
\label{subsec:routing_replay}

\paragraph{Background}
Unlike RL training in dense models, MoE models exhibit sparsely activated pathways, which bring forth distinctive challenges related to training stability. Specifically, our observations indicate that the GRPO algorithm struggles with the \textit{expert-activation volatility} characteristic of MoE architectures, undermining proper convergence during RL optimization. More concretely, following one or more gradient steps, the set of experts activated to produce an identical response can shift considerably. For instance, using the 48-layer Qwen3-30B-A3B-Base configuration, it is observed that after every RL gradient update, about 10\% of the experts involved for a fixed rollout sample under the updated policy $\pi_{\theta}$ are no longer the same as those under the previous policy $\pi_{\theta_\text{old}}$. This pattern, which intensifies as MoE models become deeper, causes severe oscillations in the token-level importance weights $w_{i,t}(\theta) = \frac{ \pi_{\theta} (y_{i,t} | x, y_{i,<t}) }{ \pi_{\theta_\text{old}} (y_{i,t} | x,y_{i,<t})}$, greatly diminishing their reliability, as elaborated in \S~\ref{sec:motivation} and~\ref{subsec:gradient}. As a result, this effect impedes standard RL training convergence.

\paragraph{Our Previous Approach} In our earlier work, we addressed this issue by adopting the \textbf{Routing Replay} training methodology. This approach involves recording the experts that are activated under $\pi_{\theta_\text{old}}$ and then reusing, or ``replaying,'' these particular routing paths under $\pi_{\theta}$ during the calculation of the importance ratios, $w_{i,t}(\theta) = \frac{ \pi_{\theta} (y_{i,t} | x, y_{i,<t}) }{ \pi_{\theta_\text{old}} (y_{i,t} | x,y_{i,<t})}$. As a result, both $\pi_{\theta} (y_{i,t} | x, y_{i,<t})$ and $\pi_{\theta_\text{old}} (y_{i,t} | x, y_{i,<t})$ are evaluated using the same expert configuration for each token $y_{i,t}$. This setup facilitates stable computation of token-level importance weights and promotes consistent optimization of the engaged expert networks throughout successive gradient steps. Figure~\ref{fig:routing_replay} illustrates that Routing Replay is crucial for ensuring reliable convergence during GRPO training of MoE models.

\paragraph{Benefit of GSPO} While Routing Replay allows GRPO training of MoE models to reach convergence, its method of reutilizing routing modes results in higher memory and communication costs and may artificially restrict the effective capacity of the MoE model. On the other hand, as depicted in Figure~\ref{fig:results}, GSPO does away with the reliance on Routing Replay by directly and reliably computing the importance ratios $s_i(\theta)$ in the standard manner, achieving stable optimization and regular convergence. A crucial observation is that GSPO exclusively depends on sequence-level likelihoods (that is, $\pi_{\theta} (y_i | x)$), rather than being impacted by token-level likelihoods (such as $\pi_{\theta} (y_{i,t} | x, y_{i,<t})$). As the MoE model consistently preserves its fundamental language modeling strength, the overall sequence likelihood remains relatively stable. To conclude, GSPO effectively addresses the volatility associated with expert activation in MoE models, making laborious techniques like Routing Replay unnecessary. This approach not only streamlines and stabilizes the training procedure but also ensures the model can fully exploit its complete representational capacity without imposed limitations.

\subsection{Benefit of GSPO for RL Infrastructure}

Due to the variations in numerical precision between engines used for training (such as Megatron) and those used for inference (like SGLang and vLLM), it is common practice to recalculate the likelihoods of generated responses under the previous policy $\pi_{\theta_\text{old}}$ using the training engine. Nevertheless, since GSPO relies on sequence-level likelihoods for its optimization process—rather than on the more sensitive token-level likelihoods—it is inherently more robust to issues stemming from precision mismatches. Consequently, GSPO allows for the direct utilization of the likelihood estimates provided by the inference engine in the optimization workflow, obviating the necessity for likelihood recomputation with the training engine. This approach proves particularly advantageous in contexts such as partial rollout, multi-turn reinforcement learning, and disaggregated frameworks that separate training and inference.

\section{Conclusion}

We introduce Group Sequence Policy Optimization (GSPO), a novel reinforcement learning technique tailored for training large language models. GSPO utilizes the core idea of importance sampling by formulating importance ratios derived from sequence probabilities and applying operations such as sequence-level clipping, rewarding, and policy optimization. When benchmarked against GRPO, GSPO achieves markedly improved training stability, greater efficiency, and enhanced performance. It proves especially effective for large-scale reinforcement learning applied to MoE models, serving as a crucial driver behind the significant progress observed in recent Qwen3 models. Leveraging GSPO as a robust and scalable approach, we are committed to advancing reinforcement learning at scale and anticipate foundational breakthroughs in AI intelligence.

\end{document}